% p1-07-mdl-bayesian

\section{P1 §7. MDL and Bayesian Evidence Framework}

7. MDL and Bayesian Evidence Framework

    Figure (fig-p1-ch7-mdl). MDL / Bayes triptych --- Two-Part Code / Bayes Factor balance / Fisher
                                              geometry.

  7.1 Model Comparison Setup

  The compressibility question is formalized as a Bayesian model comparison between:

  - M0: structureless log-uniform baseline model (null).
  - M1: structured symbolic grammar model G$\varphi$ with hypothesis class S(C).

  Both models are evaluated on the frozen target dataset T = {Xi}i=1N. The MDL
  framework (Barron, Rissanen \& Yu, 1998; Grünwald, 2007) provides the primary
  decision criterion; Bayesian evidence provides a complementary and asymptotically
  equivalent measure.

  7.2 Bayesian Evidence
  The Bayesian evidence (marginal likelihood) for each model is

  P(T mid M) = $\in$t P(T mid theta, M) P(theta mid M) dtheta. 17

  For the discrete grammar model M1, a Gibbs-type likelihood kernel with inverse
  temperature beta > 0 is adopted:

  P(T mid M1) = prodi=1N (sumF $\in$ S(C) exp[-beta d(F, Xi)])/(|S(C)|), 18

  where d(F, Xi) = |F - Xi|/|Xi| is the normalized approximation error. For the null model
  M0, the log-uniform prior gives

  P(Xi mid M0) = (1)/(log(Xmax / Xmin)). 19

  7.3 Bayes Factor and Its Calibration
  The Bayes factor comparing M1 to M0 is

  B := (P(T mid M1))/(P(T mid M0)). 20

  Standard calibration of log10 B (Grünwald, 2007):

  Crucially, B as defined in (20) does not penalize for the size of S(C). The MDL
  framework (Section 7.4) introduces the necessary penalization, without which the
  Bayes factor provides a misleadingly optimistic assessment.

  7.4 Minimum Description Length Framework
  Following Barron, Rissanen \& Yu (1998) and Grünwald (2007), the total description
  length of model M applied to data T is

  MDL(M) = L(M) + L(T mid M), 21

  where L(M) is the code length for the model (complexity cost) and L(T mid M) is the
  conditional encoding cost of the data given the model (fit cost).

  Model complexity term. For M1, the model description must specify: the basis A, the
  complexity bound C, and the selected expression F^*(theta). The code length is

  L(M1) approx log2 |S(C)| + log2 C. 22

  Data fit term. Under the Gibbs likelihood kernel,

  L(T mid M1) = -sumi=1N log2(sumF $\in$ S(C) exp[-beta d(F, Xi)]) + N log2 |S(C)|. 23

  Null model description length. Under M0,

  L(T mid M0) = N log2(log(Xmax / Xmin)), 24

  with negligible model complexity L(M0) approx 0 (no free parameters).

  Compression gain. The MDL compression advantage of the structured grammar is

  DeltaMDL := MDL(M0) - MDL(M1). 25

  A positive DeltaMDL indicates a coding advantage net of description cost. A negative or
  zero DeltaMDL provides no evidence of compressibility.

  7.5 Sensitivity to the Likelihood Temperature
  The parameter beta controls the sharpness of the approximation match. Results are
  reported for a range of beta $\in$ {10, 100, 1000} in relative-error units, and robustness of
  DeltaMDL > 0 is verified across this range. A compression advantage claim requires
  DeltaMDL > 0 for all beta values in the range and all three null-model comparisons
  simultaneously.

  7.6 Information-Geometric Interpretation
  From an information-geometry perspective, the symbolic grammar induces a
  non-Euclidean statistical manifold over R^+, with geodesic distances governed by the
  logarithmic embedding and the discrete lattice anisotropy of S(C). The observed
  compression advantage has the heuristic interpretation

  KG$\varphi$(X) $\leq$ Kunstructured(X) - DeltaMDL 26

  for well-approximated constants X, where K denotes Kolmogorov complexity in the
  sense of Chaitin (1987). The MDL quantity DeltaMDL provides a computable lower
  bound on the potential compression improvement; the inequality (26) is heuristic and
  not a theorem.

\subsection{References}

- Vasilev, D. \textit{Flos Aureus Monograph: Trinity Constants}. DOI \href{https://zenodo.org/doi/10.5281/zenodo.19227877}{10.5281/zenodo.19227877}.
- Pellis, S. "The Fine-Structure Constant and the Golden Ratio." \href{https://vixra.org/pdf/2110.0117v4.pdf}{viXra 2110.0117 v4}.
- Sherbon, M. A. "Fundamental Nature of the Fine-Structure Constant." \href{https://hal.science/hal-02341850/document}{HAL hal-02341850}.
